\documentclass[11pt,a4paper]{article}
\usepackage[T1]{fontenc}
\usepackage{lmodern}
\usepackage[margin=24mm,headheight=14pt]{geometry}
\usepackage{amsmath,amssymb,mathtools,bm}
\usepackage{booktabs,array,enumitem,microtype}
\usepackage[dvipsnames]{xcolor}
\usepackage{hyperref,fancyhdr}
\hypersetup{colorlinks=true,linkcolor=MidnightBlue,urlcolor=MidnightBlue,
 citecolor=MidnightBlue,pdftitle={Hierarchical cEBMF: capacity and paper positioning},
 pdfauthor={Research assessment for the cEBMF project}}
\pagestyle{fancy}\fancyhf{}
\fancyhead[L]{\small Hierarchical cEBMF}
\fancyhead[R]{\small Capacity and paper positioning}
\fancyfoot[C]{\thepage}
\setlength{\parindent}{0pt}
\setlength{\parskip}{5pt}
\setlist{itemsep=3pt,topsep=4pt}
\newcommand{\E}{\mathbb E}
\newcommand{\R}{\mathbb R}
\newcommand{\N}{\mathcal N}
\newcommand{\Var}{\operatorname{Var}}
\newcommand{\Cov}{\operatorname{Cov}}
\newcommand{\pa}{\operatorname{pa}}
\newcommand{\ch}{\operatorname{ch}}
\newcommand{\sigmoid}{\operatorname{sigmoid}}
\newcommand{\KL}{\operatorname{KL}}
\newcommand{\ind}{\mathbf 1}
\newcolumntype{P}[1]{>{\raggedright\arraybackslash}p{#1}}
\numberwithin{equation}{section}
\title{Hierarchical cEBMF for multi-omics\\[5pt]
\large Model capacity, analytical benchmarks, and a defensible paper contribution}
\author{Research assessment for the cEBMF project}
\date{Literature checked: 14 September 2026}
\begin{document}
\maketitle

\textbf{Assessment.} This is a plausible basis for a methods paper if the
contribution is a coherent, useful, and computationally effective solver for
\emph{learned conditional distributions between molecular programs}.
Hierarchical multi-omics, empirical Bayes integration, nonlinear latent
relationships, and interpretability each have substantial prior art. None,
on its own, is a defensible novelty claim.

The most promising combination is:
\begin{quote}
Separate additive factorizations for each modality, connected by flexible
conditional loading priors, with joint posterior inference that accounts for
uncertainty in both upstream and downstream programs.
\end{quote}
An EBNM-based computational construction, usable across the learned prior
registry, could make this statistically and practically distinctive.
It would need to work beyond the particular CGB--HMM toy example.

\textbf{Relative to your existing work.} EBMF already reduces factor
updates to normal-means problems, and cEBMF already supplies flexible
covariate-moderated priors \cite{ebmf,cebmf}. The additional contribution
would be treating the informative covariates as uncertain latent programs
inside a joint model, with the resulting feedback and parameter updates.
Prior flexibility or EBNM modularity alone would not be new.

\textbf{Closest comparisons.}
MOFA+ and JIVE are essential reference points, but a contemporary paper also
needs EV-BIDIFAC for empirical Bayes linked matrix decomposition,
CAVACHON for a hierarchical multi-omics generative model, and HALO for
relationships between RNA and ATAC representations
\cite{mofa,jive,evbidifac,cavachon,halo}.
SIFA, MuVI, SOFA, and MEFISTO cover related uses of external information
\cite{sifa,muvi,sofa,mefisto}.

\textbf{What this document establishes.}
We compare model structure, derive exact results in controlled examples,
identify the actual inference contribution, and specify an empirical
benchmark. Numerical values in the analytical examples are population
calculations, \emph{not fitted-method performance}. No competing methods or
new hierarchical solver were run for this report.

\textbf{Scope.} The proposed model below is the paired Gaussian observation
model developed in the companion
\href{run:atac_rna_joint_inference.pdf}{joint-inference derivation}.
It is a proposal, not a statement that the current package implements full
joint inference. The literature review is targeted rather than exhaustive;
publication potential is a judgment, not a guarantee of novelty or acceptance.

\clearpage
\section{The proposed model and the comparison target}

For paired samples, let
\[
Y^A\in\R^{N\times P_A},\qquad Y^R\in\R^{N\times P_R}.
\]
Use separate loading matrices \(A\in\R^{N\times K_A}\) and
\(R\in\R^{N\times K_R}\), with feature factors
\(F^m\in\R^{P_m\times K_m}\):
\begin{align}
Y^A_{ip}\mid A,F^A&\sim
 \N\!\left(\sum_{k=1}^{K_A}A_{ik}F^A_{pk},(\tau^A_{ip})^{-1}\right),\\
Y^R_{ip}\mid R,F^R&\sim
 \N\!\left(\sum_{j=1}^{K_R}R_{ij}F^R_{pj},(\tau^R_{ip})^{-1}\right).
\end{align}
Only observed entries enter the likelihood. Samples are conditionally
independent under the following row-wise prior:
\begin{equation}\label{eq:dag}
p(A,R\mid\theta)=
\prod_{i=1}^{N}
\left\{\prod_{k=1}^{K_A}
g^A_k(A_{ik}\mid A_{i,<k};\theta^A_k)\right\}
\left\{\prod_{j=1}^{K_R}
g^R_j(R_{ij}\mid R_{i,<j},A_i;\theta^R_j)\right\}.
\end{equation}
Each \(g\) is normalized; a sparse subset of the displayed parents can be
used. Fixed observed covariates may also be included. The feature-factor
prior \(p(F^A,F^R\mid\psi)\) is a separate design choice, including independent
shrinkage or an HMM along a meaningful feature ordering.

\textbf{Interpretation.} The same ATAC state can change an RNA program's
activation probability, amplitude distribution, or uncertainty. Different
RNA programs need not reuse identical ATAC scores. The model also permits
dependencies among programs within each modality.

The generative ordering \(A\to R\) does not restrict posterior inference to
that direction. For fixed parameters,
\[
p(A,R,F^A,F^R\mid Y^A,Y^R)
\ \propto\ p(Y^A\mid A,F^A)\,p(Y^R\mid R,F^R)\,
p(A,R)\,p(F^A,F^R).
\]
RNA can therefore update ATAC. This property is shared by coherent joint
latent-variable models in general; it is not unique to cEBMF.

\textbf{Four different meanings of capacity.}
\begin{enumerate}
\item \emph{Reconstruction capacity}: which finite mean matrices can be fitted?
\item \emph{Distributional capacity}: which joint laws and conditional
relationships between programs can be represented?
\item \emph{Statistical efficiency}: how much paired data is needed to learn
those relationships and improve prediction?
\item \emph{Inference capacity}: can the algorithm approximate the relevant
posterior at useful scale and quantify uncertainty?
\end{enumerate}
Greater distributional flexibility does not imply better finite-sample
performance or easier inference.

\clearpage
\section{Comparison with factorization methods}

The following are structural comparisons of the cited formulations.
An absent mechanism means it is not native to that formulation, not that
the method could never be extended to include it.

\begingroup\small
\renewcommand{\arraystretch}{1.13}
\begin{tabular}{@{}P{.19\textwidth}P{.37\textwidth}P{.38\textwidth}@{}}
\toprule
Method & Native structure and strength & Relation to the proposal\\
\midrule
JIVE (2013); AJIVE (2018)
&
Low-rank joint and individual sample subspaces; AJIVE estimates shared
structure using angles and perturbation arguments.
&
A linear shared-subspace description, without a learned nonlinear
conditional law between modality-specific scores in the original models
\cite{jive,ajive}.\\[5pt]
MOFA+ (2020)
&
Shared sample scores, view-specific feature weights, sparsity and ARD;
factors can be shared, view-specific, or active in subsets. Supports groups,
missing data and several observation likelihoods.
&
Already jointly borrows information and provides probabilistic inference.
Its standard score-prior construction does not provide the proposed
autoregressive conditional loading distributions \cite{mofa}.\\[5pt]
SLIDE (2019); BASS (2016)
&
Partially shared components across subsets of views. BASS uses Bayesian
structured shrinkage; SLIDE uses structured sparsity for the decomposition.
&
Shared/private structure and sparsity are existing capabilities, not
distinctive benefits of a loading hierarchy \cite{slide,bass}.\\[5pt]
BIDIFAC+ (2022); EV-BIDIFAC (2024)
&
Low-rank modules supported on subsets of linked row and column groups.
EV-BIDIFAC uses empirical variational Bayes shrinkage and supports
imputation.
&
A direct baseline for an empirical Bayes integration claim. The proposed
distinction is learning conditional program distributions, rather than
only the sharing pattern and shrinkage of modules
\cite{bidifac,evbidifac}.\\[5pt]
ProJIVE (2026)
&
Probabilistic joint and individual factorization, estimated by maximum
likelihood using EM.
&
A current probabilistic JIVE comparator; making JIVE probabilistic
would not be a new contribution \cite{projive}.\\
\bottomrule
\end{tabular}
\endgroup

\subsection*{A fair reading of MOFA+ and JIVE}
MOFA+ does \emph{not} require every factor to affect every modality.
JIVE explicitly includes individual variation. Therefore, allowing
modality-specific factors is insufficient positioning.

For orientation, the shared-score Gaussian predictor is
\(Y^m=ZW_m^\top+E^m\). A shared/private linear formulation writes
\(Y^m=SU_m^\top+V_m B_m^\top+E^m\).
These schematic equations describe matrix structure, not every prior,
constraint, or likelihood in the methods above.

The useful new question is whether a nonlinear, uncertain relation between
distinct programs can be learned more efficiently and interpreted more
directly through Equation~\eqref{eq:dag}. It is not whether existing models
can represent any view-specific signal at all.

\clearpage
\section{Hierarchical and information-informed comparators}

\begingroup\small
\renewcommand{\arraystretch}{1.12}
\begin{tabular}{@{}P{.19\textwidth}P{.37\textwidth}P{.38\textwidth}@{}}
\toprule
Method & Native structure and strength & Consequence for the paper\\
\midrule
CAVACHON (2024 preprint; earlier 2022 workshop)
&
Hierarchical VAE based on a supplied modality DAG; Gaussian-mixture
innovation latents feed parent-aware decoders.
&
Close prior art for biological hierarchy and conditional independence.
A factor-level conditional-prior model must be distinguished from this
decoder hierarchy \cite{cavachon}.\\[5pt]
HALO (2025)
&
RNA/ATAC representations linked through learned mappings and causal
constraint losses, with an interpretable nonlinear additive decoder.
&
Neither hierarchical RNA--ATAC modeling nor interpretable nonlinear
representations are new. Its published generative-prior and loss
construction differs from Equation~\eqref{eq:dag} \cite{halo}.\\[5pt]
SIFA (2017)
&
Joint and individual factors partly informed by observed covariates
through nonparametric functions.
&
Covariate-informed nonlinear factor relationships already exist.
Latent parents, their uncertainty, and feedback are the relevant
distinction here \cite{sifa}.\\[5pt]
MuVI (2023); SOFA (2026)
&
MuVI uses feature-set information through structured shrinkage.
SOFA separates guided and unguided factors using sample-level
guide variables.
&
Different ways to incorporate prior knowledge; include when the
evaluation supplies pathways, annotations or observed drivers
\cite{muvi,sofa}.\\[5pt]
MEFISTO (2022)
&
Gaussian-process structure on sample factors over continuous
covariates such as time and space.
&
Structured factors are existing work. Its sample-axis covariance is
different from an HMM over genomic features \cite{mefisto}.\\[5pt]
MultiVI (2023); GLUE (2022)
&
MultiVI handles paired and unpaired RNA/ATAC with a multimodal VAE.
GLUE uses a feature guidance graph to link modality embeddings.
&
Relevant for single-cell integration, count data, missing modalities
and unpaired settings. Their scope is broader than the present
Gaussian paired-data proposal \cite{multivi,glue}.\\
\bottomrule
\end{tabular}
\endgroup

\subsection*{Two distinctions that matter}
CAVACHON's generative equations put independent mixture innovation
variables into a decoder that incorporates ancestors. In contrast,
Equation~\eqref{eq:dag} places dependencies in normalized distributions
of scalar program loadings, while keeping feature reconstruction linear.
This is a different parameterization and inductive bias, not proof of
strictly greater model capacity.

HALO's Gaussian latent priors and additional constraint/mapping losses
(Methods, Equations 17--24) differ from a normalized \(p(R\mid A)\).
Neither comparison establishes superiority; both require experiments.

\clearpage
\section{Analytical benchmark I: what hierarchy does not add}

\subsection{Any finite pair of fitted matrices has a shared-score representation}
Let the fitted signals be
\[
M^A=AF_A^\top,\qquad M^R=RF_R^\top.
\]
Define \(Z=[\,A\ R\,]\) and view weights
\[
W_A=[\,F_A\ 0\,],\qquad W_R=[\,0\ F_R\,].
\]
Then
\begin{equation}
M^A=ZW_A^\top,\qquad M^R=ZW_R^\top.
\end{equation}
Thus a shared-score matrix parameterization with up to \(K_A+K_R\)
columns can reproduce the two fitted means exactly.
The minimal joint rank can be smaller.

\textbf{What follows.} There is no universal in-sample reconstruction
advantage merely from giving the modalities different loading matrices.

\textbf{What does not follow.} This construction does not match the
probability law of \(Z\), the method's prior probability on the constructed
weights, its identifiability constraints, or the efficiency of learning
these weights from new data. It is not an equality of the full MOFA+
and hierarchical cEBMF model classes.

\subsection{A linear Gaussian hierarchy reduces to shared/private structure}
For a sample, use column-vector notation and suppose
\[
a\sim\N(0,\Sigma_A),\qquad
r=B a+\eta,\qquad \eta\sim\N(0,\Sigma_\eta),\quad \eta\perp a.
\]
Then
\[
y^A=F_Aa+\epsilon_A,\qquad
y^R=F_RB a+F_R\eta+\epsilon_R.
\]
The same \(a\) drives both views, and \(\eta\) is RNA-specific.
Whitening \(a\) and \(\eta\), when their covariance matrices are
positive definite, yields independent standard Gaussian scores with
transformed feature weights.

This is a shared/private linear-Gaussian factor model at the level of the
observation distribution. Method-specific orthogonality and sparsity
restrictions can still differ. Adding a linear Gaussian hierarchy alone
therefore does not create the central expressive advantage being sought.

\subsection{Flexible ordering is not a causal identification result}
For unrestricted conditional families, any joint density can be
factorized by the chain rule in any variable order. A dense DAG is then a
parameterization, not evidence that ATAC causes RNA.
Restricted families, sparse parents, biological anchors and the choice
of factor orientation make the model informative, but introduce
assumptions that must be assessed.

Factor scaling, signs, permutations, and sometimes rotations are
non-identifiable without additional restrictions. Nonnegativity alone
does not guarantee uniqueness. Graph edges between unanchored factors
should not be presented as identified regulatory effects.

\clearpage
\section{Analytical benchmark II: nonlinear dependence and shrinkage}

\subsection{A dependency invisible to affine cross-modal prediction}
Let \(A_1,A_2\) be independent \(\operatorname{Bernoulli}(1/2)\) variables
and define an RNA loading \(R=A_1\mathbin{\mathrm{xor}}A_2\):
\[
\begin{array}{ccc}
\toprule A_1&A_2&R\\ \midrule
0&0&0\\0&1&1\\1&0&1\\1&1&0\\ \bottomrule
\end{array}
\qquad
\E R=\tfrac12,\quad
\Cov(A_1,R)=\Cov(A_2,R)=0.
\]
The population least-squares predictor among affine functions
\(b_0+b_1A_1+b_2A_2\) is \(1/2\). Its mean squared error is
\(\Var(R)=1/4\). The nonlinear oracle
\[
\E(R\mid A_1,A_2)=A_1+A_2-2A_1A_2
\]
has zero error. This establishes a strict advantage over
\emph{affine prediction using these two observed inputs}.

A neural activation prior can approximate this gate. For example,
\[
\pi_\kappa(a)=\sigmoid\{\kappa(2|a_1-a_2|-1)\},\qquad
R\mid a\sim(1-\pi_\kappa(a))\delta_0+
\pi_\kappa(a)\N(1,d^2).
\]
The Boolean model is approached as \(\kappa\to\infty\) and \(d^2\to0\);
it is not identical to finite-variance CGB. A single linear-logistic
gate cannot express XOR, whereas a suitable nonlinear network can.

\textbf{Limits of the example.} These numbers are not benchmark scores
for fitted MOFA+, JIVE, or cEBMF. An expanded shared basis containing
\(A_1A_2\) represents this mean relationship with a linear readout.
Non-Gaussian latent models can also have nonlinear posterior predictors.
No impossibility claim for every MOFA+ configuration follows.

\subsection{Useful dependence can affect variance rather than mean}
Suppose
\[
R\mid A=a\sim\N(0,v(a)),\qquad T=R+\varepsilon,\quad
\varepsilon\sim\N(0,s^2)
\]
with independent noise. Then \(\E(R\mid A)=0\), so knowing \(A\)
cannot improve prediction of \(R\) from \(A\) alone under squared error.
However,
\[
\E(R\mid T,A=a)=\frac{v(a)}{v(a)+s^2}T,\qquad
\Var(R\mid T,A=a)=\frac{v(a)s^2}{v(a)+s^2}.
\]
ATAC can therefore improve RNA denoising by setting the amount of
shrinkage, even without a conditional mean association.

This motivates conditional log scores and calibration alongside
reconstruction error. A benchmark assessing only cross-modal mean
prediction would miss an important use of conditional shrinkage.

\clearpage
\section{Analytical benchmark III: can RNA improve ATAC?}

\subsection{An exact population result}
Let \(U\) be an ATAC loading, a vector of loadings, or its noiseless
signal, with finite second moment. Let \(D_A\) denote ATAC evidence and
\(D_R\) additional RNA evidence. Under one fixed joint probability law,
define
\[
\widehat U_A=\E(U\mid D_A),\qquad
\widehat U_{AR}=\E(U\mid D_A,D_R).
\]
The tower property and orthogonality of conditional expectation give
\begin{equation}\label{eq:risk}
\boxed{\;
\E\|U-\widehat U_A\|^2-\E\|U-\widehat U_{AR}\|^2
=\E\|\widehat U_{AR}-\widehat U_A\|^2\ \geq 0.
\;}
\end{equation}
To verify this, write
\(U-\widehat U_A=(U-\widehat U_{AR})+
(\widehat U_{AR}-\widehat U_A)\).
The cross term has expectation zero after conditioning on both data sets.

Thus additional RNA evidence cannot increase the \emph{integrated
Bayes risk} of the optimal ATAC posterior mean. Improvement is strict
exactly when the posterior mean changes with positive probability.
Information can change higher moments without changing the mean.
The analogous result holds with RNA and ATAC exchanged.

The result is not a guarantee for every realized sample, for an
estimated empirical Bayes prior, or for a local variational solution.
Nor does it compare two separately fitted, possibly misspecified models.
It explains why bidirectional improvement is possible under a directed
generative model.

\subsection{Where the RNA evidence enters the actual update}
For an ATAC loading \(v=A_{ik}\), hold parameters fixed. Let \(\Omega_i^A\)
be the observed ATAC features for sample \(i\). Its exact
full conditional has the form
\begin{align}
p(v\mid\mathrm{rest})
&\propto
\exp\{-\tfrac12 a_{ik}v^2+b_{ik}v\}\,
g^A_k(v\mid A_{i,<k}) \notag\\[-2pt]
&\quad\times
\prod_{\substack{h>k\\v\in\pa(A_{ih})}}
g^A_h(A_{ih}\mid A_{i,<h})
\prod_{\substack{j\\v\in\pa(R_{ij})}}
g^R_j(R_{ij}\mid R_{i,<j},A_i).
\label{eq:fullconditional}
\end{align}
Here \(a_{ik}=\sum_{p\in\Omega_i^A}\tau^A_{ip}(F^A_{pk})^2\) and
\[
b_{ik}=\sum_{p\in\Omega_i^A}\tau^A_{ip}F^A_{pk}
\left(Y^A_{ip}-\sum_{h\ne k}A_{ih}F^A_{ph}\right).
\]
RNA observations enter indirectly through the RNA loadings and factors
being inferred jointly. The final product in
Equation~\eqref{eq:fullconditional} is the feedback missing when one only
treats estimated ATAC loadings as fixed RNA covariates.

The full conditional of a directed model already conditions on children.
Introducing an independently fitted reverse conditional does not in
general produce the same joint model; the two conditionals must be compatible.

\clearpage
\section{What an EBNM-based solver would contribute}

\subsection{A general interface, with real mathematical qualifications}
For a scalar loading \(v\), a factorized variational coordinate has
\[
q^*(v)\propto
\exp\!\left\{-\tfrac12av^2+bv+
\E_{q_{\pa(v)}}\log g(v\mid\pa(v))+
\sum_{c\in\ch(v)}
\E_{q_{\setminus v}}\log g_c(c\mid\pa(c))\right\}.
\]
The coefficients \(a,b\) now come from the \emph{expected} log likelihood:
replace the realized squared factors by their second moments, and take
expectations of the residual--factor products in \(b\). All expectations
hold \(v\) fixed. This is the standard variational coordinate identity,
not a new general theorem. For spike-and-continuous priors, densities
and updates use a common dominating measure with a zero atom.

For \emph{exact full conditionals} with other latents realized, define
the ordinary EBNM posterior at fixed parents and parameters,
\[
q_0(v)\propto e^{-av^2/2+bv}g(v\mid\pa(v)),\qquad
D(v)=\sum_{c\in\ch(v)}\log g_c(c\mid\pa(c)).
\]
Then \(p(v\mid\mathrm{rest})\propto q_0(v)e^{D(v)}\).
An independent proposal \(v'\sim q_0\) has Metropolis acceptance
\(\min\{1,\exp[D(v')-D(v)]\}\).
With suitable support and ergodicity, this targets the required
conditional; its finite-run accuracy depends on mixing.
Refitting prior parameters inside this acceptance step would invalidate
the argument. In variational inference, uncertainty in parents also
requires the expected-own-prior correction detailed in the companion note.

\subsection{The learned registry does not contain interchangeable model classes}
\begingroup\small
\renewcommand{\arraystretch}{1.10}
\begin{tabular}{@{}P{.30\textwidth}P{.65\textwidth}@{}}
\toprule
Registry family & Conditional capacity relevant to the paper\\
\midrule
\texttt{cash}, \texttt{lcash}, \texttt{po\_lcash}
& Mixture weights depend on covariates; centered scale mixtures change
sparsity/variance but retain zero conditional prior mean.\\[3pt]
\texttt{cgb}, \texttt{cgb\_sharp}
& Covariates change spike/slab probability, with global Gaussian slab
parameters. Sharpening changes fitting, not a new generic hierarchy.\\[3pt]
\texttt{cgb\_sharp\_2}
& Separate positive-mean and negative-mean Gaussian slabs plus spike;
slab means do not impose truncated support.\\[3pt]
\texttt{emdn}, \texttt{spiked\_emdn}
& Covariate-dependent mixture weights, means and variances; the latter
adds a zero atom. These are the broadest scalar families here.\\[3pt]
\texttt{hmm}, \texttt{hmm\_pos}, \texttt{hmm\_neg}
& Sequence priors, currently ignoring side covariates and using unit
adjacency. They require block treatment; in this proposal they act on \(F\).\\
\bottomrule
\end{tabular}
\endgroup

For CGB with active probability \(\pi(x)\), slab mean \(\mu\) and
variance \(d^2\),
\(\E(L\mid x)=\pi(x)\mu\) and
\(\Var(L\mid x)=\pi(x)d^2+\pi(x)(1-\pi(x))\mu^2\).
This is a gate, not unrestricted regression.
The current registry's posterior moments alone are not enough for
general child messages: density evaluation, sampling or integration,
and coherent parameter learning are needed.

\clearpage
\section{An empirical benchmark that would test the claims}

These are proposed experiments, not predicted rankings.
The central design is to vary \emph{dependence type} independently of
noise level, rank, sparsity, and feature smoothness.

\begingroup\small
\renewcommand{\arraystretch}{1.13}
\begin{tabular}{@{}P{.26\textwidth}P{.35\textwidth}P{.33\textwidth}@{}}
\toprule
Scenario & Question tested & Informative comparisons\\
\midrule
Linear Gaussian shared/private factors
& Does flexibility cost accuracy or runtime when unnecessary?
& MOFA+, AJIVE/ProJIVE, EV-BIDIFAC; linear hierarchy.\\[4pt]
Nonlinear activation: AND, XOR, smooth gates
& Can conditional program structure improve denoising and prediction?
& Conditional versus independent priors; HALO, CAVACHON;
shared factors plus nonlinear regression.\\[4pt]
Variance-only dependence
& Does the method learn adaptive shrinkage and predictive dispersion?
& Conditional scale mixture versus homoscedastic/independent priors.\\[4pt]
Asymmetric SNR, reversed in a second experiment
& Can downstream evidence help the noisy upstream modality?
& Full joint versus one-way plug-in and no-child-message ablations.\\[4pt]
No dependence; shuffled pairing; wrong graph
& Can the method decline to borrow, and resist negative transfer?
& Independent fits and sparsity-controlled hierarchy.\\[4pt]
Three or more modalities with subset sharing
& Does a DAG improve on established partial-sharing structure?
& SLIDE, BASS, EV-BIDIFAC; matched total latent dimension.\\[4pt]
Block features, smooth features, and permuted features
& Is any gain due to the loading hierarchy or the HMM?
& Same loading model with and without the feature prior; smoothing controls.\\[4pt]
Realistic count noise, batch effects, missing modalities
& What is the scope beyond normalized Gaussian matrices?
& MultiVI; HALO/CAVACHON for paired hierarchy; explicit
preprocessing and likelihood controls.\\
\bottomrule
\end{tabular}
\endgroup

\subsection*{Required internal ablations}
Use the same observation model and feature prior for:
(i) independent cEBMF;
(ii) one-way ATAC-to-RNA plug-in conditioning;
(iii) one-way conditioning in the reverse direction;
(iv) the proposed joint inference;
(v) within-modality hierarchy only;
(vi) a linear conditional prior.
Also compare an independent factorization followed by a trained nonlinear
map between loadings. That tests whether the benefit requires joint
inference rather than just a flexible regression.

The current block-shaped toy data are valuable for debugging, but align
closely with the HMM prior. They do not establish general multi-omics
performance. Vary the feature prior independently of the loading model.
RNA columns need a scientifically meaningful ordering before adjacency
smoothing is interpreted biologically; unit index spacing is not genomic
distance.

\clearpage
\section{Evaluation, identifiability and computational evidence}

\textbf{Primary endpoints.} In simulations, report noiseless-signal
reconstruction error separately for ATAC and RNA, conditional prediction
error in both directions, and uncertainty calibration. Assess program
recovery after appropriate scale, sign and permutation alignment; use
subspace recovery where individual factors are not identifiable.
Log scores require predictive distributions; original JIVE does not
natively provide them.

\textbf{Three distinct holdouts.}
\begin{enumerate}
\item Hold out entries for within-sample denoising.
\item Hold out an entire modality for test samples, inferring their
loadings only from available modalities.
\item Hold out biological samples, donors or studies for transfer.
For single-cell studies, avoid treating cells from the same donor as
independent biological replications.
\end{enumerate}
The corresponding estimands are different and should not be pooled.
An entire-modality holdout on genuinely new samples is especially
informative for learned conditional distributions.

\textbf{Prevent leakage.} Train transformations, feature selection,
hyperparameters and conditional networks using training data; select
settings with validation data. Test RNA must not enter the ATAC covariates
used to predict that same held-out RNA. Parent uncertainty should be
integrated, not replaced by information from the hidden modality.
Distinguish transductive matrix completion from prediction for new samples.

\textbf{Fairness.} Tune ranks and regularization rather than fixing an
arbitrary common \(K\). A model with \(K_A\) ATAC factors and \(K_R\) RNA
factors does not have the same complexity as a shared model with
\(\max(K_A,K_R)\) factors. Report both matched-capacity and independently
tuned comparisons, network sizes, optimization budgets and restarts.
Supply the same external information when comparing mechanisms that use
it. When HMM smoothing is unique to one arm, report that advantage separately.

\textbf{Comparable scales.} Raw ELBO values from different model families
are not a common performance metric. Compare held-out prediction for the
same observed quantity and units. Gaussian losses on transformed
expression are not directly comparable to count likelihoods.
Use repeated simulation replicates with uncertainty intervals, and
held-out donors/studies for real-data uncertainty.
Clustering and enrichment can supplement, but cannot replace, prediction
and calibration.

\textbf{Computational evidence.} A residual-based Gaussian sweep has a
likelihood cost of order \(N\sum_m P_mK_m\). The loading hierarchy adds
child-prior evaluations and parameter optimization. With \(K=K_A+K_R\),
a dense ordered graph has \(K(K-1)/2\) edges, so cost can grow rapidly with
rank. Integration accuracy and network input size add further costs.
An HMM with \(S\) states has the usual \(O(P_mS^2)\) forward--backward
cost per feature-factor sequence before exploiting transition structure.
Report runtime, memory, convergence and, for Monte Carlo, effective sample
size. No scaling advantage can be claimed from the update formula alone.

\textbf{Suggested real-data scope.} Start with two independent paired
RNA--ATAC studies, preferably spanning different tissues or technologies.
Add a three-modality application only if broad multi-omics capability is
a central claim. For the first paper, a carefully delimited Gaussian
analysis of processed molecular features is legitimate; raw-count
performance needs an observation-model extension or a justified pipeline.

\clearpage
\section{What would make this a convincing paper?}

\subsection*{A defensible contribution statement}
\begin{quote}
We develop a hierarchical empirical Bayes factorization that learns
conditional distributions between modality-specific molecular programs.
The model preserves additive feature contributions and supports joint
posterior updates across a specified program graph. An EBNM-based
inference strategy accommodates multiple conditional prior families.
We evaluate when this structure improves estimation and when it causes
negative transfer.
\end{quote}
This is suggested positioning, not a claim that all of it has been
implemented or validated.

\textbf{The most substantial methodological contribution would be the
solver and its statistical behavior.} Standard full-conditionals, the
variational coordinate identity, and the conditional-expectation risk
identity are established tools. Their application is useful, but a paper
needs more than writing them down.

A strong submission would provide:
\begin{enumerate}
\item A normalized joint model, clearly separating observed covariates,
latent parents, feature structure, and learned hyperparameters.
\item An implemented common-objective inference algorithm, with mixed
spike/continuous supports, uncertain-parent expectations and child
feedback handled consistently. A small exact or accurately sampled
reference should check the approximation.
\item Evidence across at least two materially different conditional
prior families, for example CGB gates and mixture mean/variance models.
Regularization and preprocessing must have a coherent role in the target.
\item A clear account of ordering, anchors, sparsity and factor
identifiability; sensitivity to alternative orderings and incorrect
graphs should be reported.
\item Fair comparisons against strong linear/EB methods and hierarchical
neural comparators, with convincing positive examples and failures.
\end{enumerate}

\textbf{Useful theory targets.} Characterize a restricted model class for
which conditional program dependencies are identifiable, or prove
approximation/error properties for the proposed EBNM message algorithm.
An exact-coordinate ascent proof gives monotonicity only for updates
that truly maximize one well-defined objective; it does not automatically
cover Monte Carlo gradients, plug-in parents, sharpening heuristics,
changing supports, or repeated black-box prior refits.

\textbf{My judgment.} The idea is promising as an interpretable statistical
alternative for studying dependence between programs. Its clearest
potential benefit over standard shared-factor models is a targeted
conditional-distribution assumption, especially with nonlinear gates or
covariate-dependent shrinkage. Its benefit over hierarchical VAEs would
need to come from the particular additive model, prior flexibility,
inference quality, data efficiency, or practical simplicity.

A paper showing only that one alternating two-fitter example improves RNA
would be weak. A carefully validated general solver, a useful biological
question, and evidence explaining \emph{when and why} joint conditional
modeling helps could constitute a substantial contribution.

\clearpage
\begingroup\small\raggedright
\begin{thebibliography}{99}
\setlength{\itemsep}{5pt}
\bibitem{mofa}
Argelaguet et al. (2020).
MOFA+: a statistical framework for comprehensive integration of
multi-modal single-cell data. \emph{Genome Biology} 21, 111.
\url{https://doi.org/10.1186/s13059-020-02015-1}

\bibitem{jive}
Lock, Hoadley, Marron and Nobel (2013).
Joint and Individual Variation Explained (JIVE) for integrated analysis
of multiple data types. \emph{Annals of Applied Statistics} 7, 523--542.
\url{https://pmc.ncbi.nlm.nih.gov/articles/PMC3671601/}

\bibitem{ajive}
Feng, Jiang, Hannig and Marron (2018).
Angle-based joint and individual variation explained.
\url{https://arxiv.org/abs/1704.02060}

\bibitem{slide}
Gaynanova and Li (2019).
Structural learning and integrative decomposition of multi-view data.
\emph{Biometrics} 75, 1121--1132.
\url{https://arxiv.org/abs/1707.06573}

\bibitem{bass}
Zhao, Gao, Mukherjee and Engelhardt (2016).
Bayesian group factor analysis with structured sparsity.
\emph{Journal of Machine Learning Research} 17.
\url{https://www.jmlr.org/papers/volume17/14-472/14-472.pdf}

\bibitem{bidifac}
Lock et al. (2022).
Decomposition and imputation of bidimensionally linked matrices.
\url{https://pmc.ncbi.nlm.nih.gov/articles/PMC9060567/}

\bibitem{evbidifac}
Lock (2024).
Empirical Bayes Linked Matrix Decomposition.
\emph{Machine Learning} 113, 7451--7477.
\url{https://pmc.ncbi.nlm.nih.gov/articles/PMC11698509/}

\bibitem{projive}
Murden et al. (2026).
Probabilistic Joint and Individual Variation Explained (ProJIVE) for
Data Integration. \emph{Journal of Computational and Graphical Statistics};
arXiv version accessed.
\url{https://arxiv.org/abs/2603.12351}

\bibitem{cavachon}
Hsieh et al. (2024).
CAVACHON: a hierarchical variational autoencoder to integrate multi-modal
single-cell data. Preprint, arXiv:2405.18655v1.
Earlier related paper presented at NeurIPS LMRL Workshop 2022.
\url{https://arxiv.org/html/2405.18655v1}

\bibitem{halo}
Mao et al. (2025).
HALO: hierarchical causal modeling for single cell multi-omics data.
\emph{Nature Communications} 16, 8892.
\url{https://www.nature.com/articles/s41467-025-63921-1}

\bibitem{sifa}
Li and Jung (2017).
Incorporating Covariates into Integrated Factor Analysis of Multi-View
Data. \emph{Biometrics} 73, 1433--1442.
\url{https://arxiv.org/abs/1703.05794}

\bibitem{muvi}
Qoku and Buettner (2023).
Encoding Domain Knowledge in Multi-view Latent Variable Models:
A Bayesian Approach with Structured Sparsity. \emph{AISTATS}, PMLR 206.
\url{https://proceedings.mlr.press/v206/qoku23a.html}

\bibitem{sofa}
Capraz et al. (2026).
Semi-supervised Omics Factor Analysis (SOFA) disentangles known and latent
sources of variation in multi-omic data.
\emph{Nature Communications} 17, 8725. Published 20 August 2026.
\url{https://www.nature.com/articles/s41467-026-74694-6}

\bibitem{mefisto}
Velten et al. (2022).
Identifying temporal and spatial patterns of variation from multimodal
data using MEFISTO. \emph{Nature Methods} 19, 179--186.
\url{https://pmc.ncbi.nlm.nih.gov/articles/PMC8828471/}

\bibitem{multivi}
Ashuach et al. (2023).
MultiVI: deep generative model for the integration of multimodal data.
\emph{Nature Methods}.
\url{https://www.nature.com/articles/s41592-023-01909-9}

\bibitem{glue}
Cao and Gao (2022).
Multi-omics single-cell data integration and regulatory inference with
graph-linked embedding. \emph{Nature Biotechnology}.
\url{https://pmc.ncbi.nlm.nih.gov/articles/PMC9546775/}

\bibitem{ebmf}
Wang and Stephens (2021).
Empirical Bayes Matrix Factorization.
\emph{Journal of Machine Learning Research} 22(120), 1--40.
\url{https://jmlr.org/papers/v22/20-589.html}

\bibitem{cebmf}
Denault et al. (2025).
Covariate-moderated Empirical Bayes Matrix Factorization.
Preprint, arXiv:2505.11639v2.
\url{https://arxiv.org/html/2505.11639v2}
\end{thebibliography}
\endgroup
\end{document}
